\subsection{Model Hierarchy and Comparison Boundary}
\label{subsec:model-hierarchy-comparison-boundary}

Reduced hemodynamic models form a hierarchy of modeling roles rather than a
scalar accuracy ranking. The governing equations above start from a
three-dimensional continuum problem and then introduce reduced descriptions, so
this subsection reads the hierarchy in descending spatial dimension. Dimension,
wall treatment, rheology, boundary data, coupling laws, output operators, and
numerical realization remain separate modeling choices. This report selects the
1D area--flow model as the reduced forward map; 3D, multiscale, 2D, and 0D
models supply comparison, coupling, diagnostic, and boundary-condition
vocabulary.

\begin{figure}[!htb]
    \centering
    \captionsetup{aboveskip=2pt,belowskip=2pt}
    \figtikz{model-hierarchy}
    \caption{Model hierarchy for reading comparison claims, ordered from the
    three-dimensional continuum tier toward reduced and lumped closures. The
    horizontal direction is a dimensional-reduction and closure-abstraction
    axis, not a validation ranking. This report selects a 1D reduced forward
    model; 3D fixed-wall, resolved FSI, clinical, and independent-solver
    comparisons require separately matched assumptions, inputs, output
    operators, and discrepancy metrics before they can support evidence claims.}
    \label{fig:model-hierarchy-overview}
\end{figure}

\begin{definition}[Wall and coupling closure axes]
    \label{def:wall-fsi-closure-catalog}
    The wall descriptor $\mathcal W$ records how the wall or wall-reduced
    mechanics are closed. The coupling descriptor $\mathcal C$ records whether
    the single vessel is isolated or coupled to another model tier. Both are
    independent of the rheology descriptor $\mathcal R$: $\mathcal R$ changes
    the stress or effective-viscosity law, $\mathcal W$ changes the wall
    mechanics, and $\mathcal C$ changes the interface or network contract.
    \begin{center}
    \footnotesize
    \begin{tabularx}{\textwidth}{@{}
        >{\raggedright\arraybackslash}p{0.17\textwidth}
        >{\raggedright\arraybackslash}p{0.18\textwidth}
        >{\raggedright\arraybackslash}X
        >{\raggedright\arraybackslash}X
        @{}}
        \toprule
        \textbf{Axis}
            & \textbf{Descriptor}
            & \textbf{Information retained}
            & \textbf{Reading in this report} \\
        \midrule
        Wall
            & $\mathcal W_{\mathrm{rigid}}$
            & Fixed or prescribed wall geometry; no structural state is solved
            & A fixed-wall continuum or rigid-tube reduced baseline. \\
        Wall
            & $\mathcal W_{\mathrm{elastic1D}}$
            & Dynamic area $A(z,t)$, gauge pressure $p_{\mathrm g}(A,z,t)$,
            reference area $A_0(z)$, wall stiffness $\beta(z)$, and external
            pressure
            & A pressure--area wall-law family. The selected 1D model below
            uses its Canic--Koiter member; no separate structural PDE is
            solved. \\
        Wall
            & $\mathcal W_{\mathrm{resolvedFSI}}$
            & Fluid velocity and pressure plus wall displacement, wall velocity,
            structural stress, and interface conditions
            & A full moving-wall FSI PDE system, used here only as comparison
            vocabulary. \\
        Coupling
            & $\mathcal C_{\mathrm{isolated}}$
            & Inlet and outlet data close one vessel
            & The baseline numerical study. \\
        Coupling
            & $\mathcal C_{\mathrm{network}}$ or $\mathcal C_{\mathrm{3D/1D}}$
            & Interface pressure, flow, traction, impedance, or characteristic
            data connecting 0D, 1D, and 3D tiers
            & A separate boundary/interface contract, not a wall law. \\
        \bottomrule
    \end{tabularx}
    \end{center}
    This separation follows the standard distinction between reduced compliant
    vessel models, geometrical multiscale coupling, and resolved FSI models
    \cite[Sec.~2, pp.~253--257]{FormaggiaEtAl2003OneDimensionalBloodFlow}
    \cite{QuarteroniVeneziani2003GeometricalMultiscale}
    \cite{FormaggiaEtAl2007FSICoupling}.
\end{definition}

\subsubsection{3D Fixed-Wall and Resolved-FSI Models}
\label{subsubsec:3d-models}

A fixed-wall 3D hemodynamic model solves for the velocity field
$\vu(t,\vx)$, pressure field $p(t,\vx)$, and stress-derived outputs on a
three-dimensional lumen with fixed or prescribed wall geometry
\cite{QuarteroniFormaggia2004CardiovascularModeling}. It can supply resolved
velocity, pressure, and wall traction. A WSS output additionally requires the
stress tensor, wall normal, wall location, tangent convention, and output
operator to be declared
\cite{JohnEtAl2017WallShearStressVectorForm}. These quantities are fields on
the resolved lumen and describe the vessel wall's interaction with the fluid
domain, not the 1D state variables $A$, $Q$, and $p_{\mathrm g}$.

Resolved FSI changes the wall descriptor. A fixed-wall model belongs to the
$\mathcal W_{\mathrm{rigid}}$ or prescribed-wall reading in
Definition~\ref{def:wall-fsi-closure-catalog}; a resolved-FSI model uses
$\mathcal W_{\mathrm{resolvedFSI}}$ and adds wall displacement, wall velocity,
structural stress, and fluid--structure interface conditions. The elastic 1D
wall law $\mathcal W_{\mathrm{elastic1D}}$ is a reduced wall closure, not a
structural FSI solve
\cite{FormaggiaEtAl2007FSICoupling}.

The 3D/FSI tier can represent eccentricity, curvature, branches, secondary
flow, separation, and recirculation more directly than a 1D area--flow model.
That role does not make it assumption-free. A valid comparison requires matched
geometry, rheology descriptor $\mathcal R$, wall closure $\mathcal W$, boundary
data, mesh and time resolution, output operators, and discrepancy metrics. A
velocity-slab diagnostic, a traction-based WSS comparison, and a pressure-ratio
comparison are different observation problems.

The 3D/FSI tier supplies the reference vocabulary for resolved velocity,
pressure, traction, and wall motion, but it does not automatically validate a
reduced model without matched comparison data. It is not the selected broad
parametric generator in this report because the setup and computational burden
are high, not because resolved 3D or FSI models are unimportant.

\subsubsection{Multiscale 3D--1D--0D Coupling}
\label{subsubsec:multiscale-models}

A multiscale model is a coupled forward problem that combines tiers of the
hierarchy. A common arrangement resolves a local 3D stenosis region, connects it
to a surrounding 1D distributed vessel network, and closes terminal beds with
0D Windkessel or related outlet models. The purpose is to keep local resolved
physics where it matters while representing wave propagation and distal
impedance without making the entire circulation a 3D domain
\cite{QuarteroniVeneziani2003GeometricalMultiscale}.

The interface laws are part of the model. A 3D--1D interface may impose
conservation of flow together with pressure, mean-pressure, normal-traction, or
characteristic compatibility. A 1D--0D interface passes terminal flow into an
ODE load and returns a pressure or impedance relation to the 1D boundary.
Resolved-FSI coupling additionally requires compatible wall kinematics,
structural stresses, and interface work. These choices affect stability and
interpretation, and they must be recorded with the numerical coupling convention
\cite{FormaggiaEtAl2007FSICoupling}.

A multiscale model is not a new fidelity tier; it is a coupled forward problem
whose interface laws are part of the model definition. Coupling more components
does not automatically make a comparison more correct. The boundary data,
interface variables, time coupling, output maps, and discrepancy metrics decide
what claim the computation can support. In this report, multiscale coupling
provides comparison and boundary-condition vocabulary; it is not the active
baseline single-vessel solve unless explicitly declared.

\subsubsection{2D Intermediate Resolved-Flow Models}
\label{subsubsec:2d-models}

A 2D model is treated here as a conceptual intermediate resolved-flow tier,
usually posed as an axisymmetric or planar flow problem. It retains a spatially
resolved velocity and pressure field on an idealized reduced domain rather than
only the section-averaged state $(A,Q,p_{\mathrm g})$. Within those
idealizations, it can be used to inspect radial velocity structure, throat
acceleration, simplified recirculation, and wall-shear trends that a 1D
area--flow model cannot resolve.

The same idealization also fixes the boundary of the evidence. Axisymmetric and
planar models suppress eccentric plaques, curvature, torsion, branch geometry,
non-axisymmetric secondary flow, and fully patient-specific WSS\@. Their wall
outputs depend on the declared stress tensor, wall location, normal, tangent,
and output operator, just as in the shear convention of
Convention~\ref{conv:appendix-shear-diagnostic}. A 2D result is therefore not
interchangeable with matched 3D or resolved-FSI evidence. The broad continuum
and dimension-reduction background for this reading is supplied by
\cite{GaldiEtAl2008HemodynamicalFlows,KopplHelmig2023DimensionReducedModeling}.

The 2D tier is a diagnostic bridge: useful for controlled mechanism checks, but
too idealized to serve as clinical or resolved-3D evidence. In this report it
can help interpret whether a reduced trend is plausibly tied to throat
acceleration or simplified recirculation, but it is not the selected forward
generator and does not validate the 1D map by itself.

\subsubsection{1D Distributed Area--Flow Models}
\label{subsubsec:1d-models}

The selected 1D reduction is defined in
Section~\ref{subsubsec:cross-sectional-reduction-selected-1d-model}. There the
retained fields $A(z,t)$, $Q(z,t)$, $\bar u(z,t)=Q/A$, and
$p_{\mathrm g}(z,t)$, the wall and coupling closure catalog, the elastic-wall
area--flow balance law, the stenosis reference geometry, the conservative
flux/source notation, and the single-vessel boundary data are recorded in
Definitions~\ref{def:section-averaged-1d-fields}--\ref{def:1d-single-vessel-boundary-data}
\cite{FormaggiaEtAl2003OneDimensionalBloodFlow,KopplHelmig2023DimensionReducedModeling}.
This subsection uses that machinery only to identify the hierarchy role of the
1D tier.

In the hierarchy, the key point is that the 1D state is distributed but reduced.
It retains axial area, flow, mean velocity, and gauge pressure; it does not
retain the full three-dimensional velocity field, pressure field, secondary
flow, separation, recirculation, or wall traction. Reduced shear quantities
therefore inherit the declared velocity-profile, wall-friction, and rheology
closures. They are not traction-based WSS unless a separate observation map
declares how a reduced output is compared with a resolved wall-stress quantity.

The closure contract remains part of the model definition. A 1D case must
declare the momentum-flux/profile convention, friction law, rheology descriptor
$\mathcal R$, wall descriptor $\mathcal W$, boundary data, numerical
realization, and output operators. Stenosis enters through the reference
geometry $R_0(z)$, $A_0(z)$, and $s(z)$, but resolved throat flow does not
appear automatically; variable-radius corrections such as the Canic-style terms
are additional closure choices
\cite{CanicEtAl2024Extended1DStenosis}.
The current implementation records this axis explicitly through the
\texttt{canic-extended-1d} and \texttt{classical-1d-no-slip} forward-model
descriptors, the latter corresponding to the wall no-slip baseline in
Definition~\ref{def:classical-no-slip-1d-baseline}.

The 1D model is the selected reduced forward map because it retains distributed
pressure--flow structure at low computational cost, while its outputs remain
reduced quantities with declared observation maps. In this report it is the
selected single-vessel generator for the numerical study, not a
clinical-validation claim, an independent-solver parity claim, or a
resolved-3D-accuracy claim.

\subsubsection{0D Outlet and Network Models}
\label{subsubsec:0d-models}

A 0D model is a lumped ODE or circuit model for pressures, flows, volumes, and
loads. It retains compartment pressures, branch flows, resistances,
compliances, and, when the circuit includes inertial effects, inertances. The
model preserves a time-dependent pressure--flow relation for parts of the
circulation that are not spatially resolved. Windkessel and RCR models are the
canonical terminal examples: they supply an outlet impedance or downstream load
for a local 1D or 3D vessel solve
\cite{ShiEtAl2011BloodFlowModelReview,KopplHelmig2023DimensionReducedModeling}.

A 0D model closes the circulation seen by the local vessel model; it does not
resolve the local stenosis. The spatial coordinate has been integrated out, so
the model has no local stenosis geometry, no axial pressure-drop field, no
throat acceleration profile, and no recirculation region. It also has no
resolved wall-shear stress or wall traction, because it does not carry a lumen
surface, a stress tensor on that surface, or a wall-normal output operator.

The 0D tier is therefore useful boundary and network context for this report. It
can prescribe an outlet load, organize distal-bed assumptions, or provide a
lumped circulation closure. It is not the selected local stenosis-resolution
tier and cannot replace the distributed area--flow state used in the numerical
study in this report.


A comparison claim requires matched assumptions and declared output operators:
the continuum setting, geometry, rheology descriptor $\mathcal R$, wall
descriptor $\mathcal W$, coupling descriptor $\mathcal C$, boundary data,
numerical settings, observation maps, and discrepancy metrics must all be part
of the comparison record. This report treats the selected 1D area--flow model
as the forward map for the numerical study. Resolved-3D/FSI agreement, clinical
calibration, and independent-solver parity are later evidence standards, not
assumptions built into the model hierarchy.
